\section*{Chapitre \ref{ch:g10:exercise-and-energy} --- L'effort et les besoins énergétiques du corps}

\begin{solution}{exo:g10:exercise-and-energy:1}
L'énergie dépensée au repos complet pour garder le corps en vie ---
cœur, cerveau, température, entretien des \omterm{def:g10:cells-common-unit:cell}{cellules}~: environ
\qty{80}{W}, soit quelque \qty{7000}{kJ} par jour.
\end{solution}

\begin{solution}{exo:g10:exercise-and-energy:2}
$2.0 \times 20 = \qty{40}{kJ}$ par minute, soit $40000/60 \approx
\qty{670}{W}$.
\end{solution}

\begin{solution}{exo:g10:exercise-and-energy:3}
La réserve phosphatée du muscle (une dizaine de secondes)~; la
\omterm{def:g10:cell-metabolism:fermentation}{fermentation} lactique (une à deux minutes d'effort maximal)~; la
\omterm{prop:g10:cell-metabolism:respiration}{respiration cellulaire} (des heures).
\end{solution}

\begin{solution}{exo:g10:exercise-and-energy:4}
Le débit le plus élevé auquel le corps peut prélever et utiliser
l'oxygène. $3500/70 = \qty{50}{mL/min}$ par kilogramme.
\end{solution}

\begin{solution}{exo:g10:exercise-and-energy:5}
Dans les muscles, environ \qty{400}{g}, pour alimenter le muscle qui le
stocke~; dans le foie, environ \qty{100}{g}, libéré sous forme de
glucose dans le sang pour le cerveau et les autres tissus.
\end{solution}

\begin{solution}{exo:g10:exercise-and-energy:6}
\qty{2.7}{L/min}~: $2.7 \times 20 = \qty{54}{kJ}$ par minute, soit
\qty{900}{W}. Chaleur~: $900 - 200 = \qty{700}{W}$.
\end{solution}

\begin{solution}{exo:g10:exercise-and-energy:7}
Travail $mgh = 60 \times 9.8 \times 800 \approx \qty{470}{kJ}$. À 25~\%
de rendement, les muscles ont dépensé environ \qty{1900}{kJ}. Sur
\qty{7200}{s}~: environ \qty{260}{W} au-dessus du repos.
\end{solution}

\begin{solution}{exo:g10:exercise-and-energy:8}
Un 100 mètres tourne presque entièrement sur la réserve phosphatée~:
pas de dette d'oxygène, peu d'acide lactique. Un 400 mètres dure environ
une minute et tire la moitié de son énergie de la \omterm{def:g10:cell-metabolism:fermentation}{fermentation}
lactique~: l'acide lactique s'accumule (la douleur) et la respiration
qui doit l'éliminer laisse le coureur haletant plusieurs minutes après.
\end{solution}

\begin{solution}{exo:g10:exercise-and-energy:9}
Énergie glucidique $0.6 \times 3000 = \qty{1800}{kJ}$~; \omterm{def:g10:exercise-and-energy:glycogen}{glycogène}
$1800/17 \approx \qty{106}{g}$~; eau libérée environ \qty{320}{g}.
\end{solution}

\begin{solution}{exo:g10:exercise-and-energy:10}
Puissance métabolique d'environ \qty{600}{W}~: $\num{1.1e6}/600 \approx
\qty{1830}{s}$, soit environ 30 minutes.
\end{solution}

\begin{solution}{exo:g10:exercise-and-energy:11}
Les trois quarts de l'énergie libérée dans les muscles sont de la
chaleur, soit plusieurs centaines de watts à l'effort intense. La sueur
l'évacue~: évaporer de l'eau demande environ \qty{2.4}{kJ} par gramme,
et le débit sanguin vers la peau y porte la chaleur. Sans sueur, la
température du corps monterait d'un degré toutes les quelques minutes.
\end{solution}

\begin{solution}{exo:g10:exercise-and-energy:12}
En montée, la puissance nécessaire est proportionnelle à la masse
corporelle~: le sujet de \qty{60}{kg} dispose de \qty{60}{mL/min} par
kilogramme contre \num{40}, et il monte plus vite. Sur le vélo de
laboratoire, la puissance ne dépend pas de la masse, tous deux ont les
mêmes \qty{3.6}{L/min}, et tous deux peuvent soutenir les mêmes watts.
\end{solution}

\begin{solution}{exo:g10:exercise-and-energy:13}
Les \qty{0.4}{L/min} d'équivalent oxygène qui manquent sont fournis par
la \omterm{def:g10:cell-metabolism:fermentation}{fermentation} lactique dans les muscles. L'acide lactique s'accumule~;
en quelques minutes, les muscles brûlent et le coureur est contraint de
ralentir jusqu'à une allure que la respiration seule peut couvrir.
\end{solution}

\begin{solution}{exo:g10:exercise-and-energy:14}
Les salves maximales courtes tournent sur la réserve phosphatée et la
\omterm{def:g10:cell-metabolism:fermentation}{fermentation}, qui n'utilisent aucune graisse~; la graisse n'est brûlée
que par la respiration, qui met des minutes à atteindre son plein
régime et ne domine que dans l'effort modéré prolongé. L'énergie totale
de quelques salves est d'ailleurs faible devant celle d'une heure
d'effort régulier.
\end{solution}

\begin{solution}{exo:g10:exercise-and-energy:15}
Déficit $25000 - 8500 - 6000 = \qty{10500}{kJ}$, soit environ
\qty{280}{g} de graisse. Manger pendant l'étape maintient la glycémie
pour le cerveau et épargne le \omterm{def:g10:exercise-and-energy:glycogen}{glycogène} pour les cols, où la puissance
nécessaire dépasse de loin ce que la graisse seule peut fournir~; la
graisse couvre la partie régulière de la journée.
\end{solution}

\begin{solution}{pb:g10:exercise-and-energy:1}
\textbf{1.} $4.0 \times 60 \times 42.2 \approx \qty{10100}{kJ}$.

\textbf{2.} $\qty{3}{h}\,\qty{30}{min} = \qty{12600}{s}$~:
$\num{1.01e7}/12600 \approx \qty{800}{W}$.

\textbf{3.} $10100/20 \approx \qty{506}{L}$~; sur \qty{210}{min},
environ \qty{2.4}{L/min}.

\textbf{4.} $2400/60 = \qty{40}{mL/min}$ par kilogramme, soit $40/52
\approx 77\%$ de son maximum.

\textbf{5.} $(2.4 - 0.2)/2.4 \approx 92\%$.

\textbf{6.} $380 \times 17 = \qty{6460}{kJ}$.

\textbf{7.} La course coûte \qty{240}{kJ} par kilomètre~: $6460/240
\approx \qty{27}{km}$.

\textbf{8.} Consommation glucidique $0.8 \times 240 = \qty{192}{kJ/km}$~:
$6460/192 \approx \qty{34}{km}$ --- le mur.

\textbf{9.} Énergie de la graisse $0.2 \times 10100 \approx \qty{2020}{kJ}$,
soit $2020/38 \approx \qty{53}{g}$~: environ 0.6~\% de ses \qty{9}{kg}.

\textbf{10.} La graisse n'est brûlée que par la respiration et à débit
limité~: elle peut fournir en gros la moitié de la puissance qu'exige
son allure. Sans \omterm{def:g10:exercise-and-energy:glycogen}{glycogène}, les muscles ne peuvent pas fabriquer l'ATP
assez vite, et l'allure tombe.

\textbf{11.} $60 \times 3.5 = \qty{210}{g}$, soit $210 \times 17 \approx
\qty{3570}{kJ}$, ce qui couvre $3570/192 \approx \qty{19}{km}$ de
consommation glucidique.

\textbf{12.} \omterm{def:g10:chemistry-of-life:families}{Glucides} disponibles $6460 + 3570 = \qty{10030}{kJ}$~:
$10030/192 \approx \qty{52}{km}$, au-delà de l'arrivée. Oui.

\textbf{13.} $+\qty{200}{g}$ de \omterm{def:g10:exercise-and-energy:glycogen}{glycogène} $+ \qty{600}{g}$ d'eau~:
\qty{0.8}{kg}, ce qui coûte $0.8 \times 4.0 \approx \qty{3}{kJ}$ de
plus par kilomètre --- négligeable devant 240.

\textbf{14.} \omterm{def:g10:exercise-and-energy:glycogen}{Glycogène} $580 \times 17 = \qty{9860}{kJ}$ plus la
boisson~: environ \qty{13400}{kJ}~; consommation glucidique
$0.55 \times 240 =
\qty{132}{kJ/km}$~: plus de \qty{100}{km}. Le mur s'est déplacé bien
au-delà de l'arrivée, avec une large marge.

\textbf{15.} Le cerveau ne brûle que du glucose, fourni par le
\omterm{def:g10:exercise-and-energy:glycogen}{glycogène} du foie. Quand cette réserve est vide, la glycémie tombe, et
le cerveau signale l'urgence par des vertiges et une fatigue
écrasante, quoi que les muscles détiennent encore.

\textbf{16.} $3.6 \times 55 \times 42.2 \approx \qty{8360}{kJ}$~;
$8360/20 = \qty{418}{L}$ sur \qty{125}{min}~: environ \qty{3.3}{L/min}.

\textbf{17.} $3340/55 \approx \qty{61}{mL/min}$ par kilogramme, soit
$61/80 \approx 76\%$ --- la même fraction que Nadia. C'est le plafond
qui diffère, non la fraction utilisée.

\textbf{18.} L'économie de course est le coût par kilogramme et par
kilomètre. Pour un coureur de \qty{55}{kg},
$(4.0 - 3.6) \times 55 \times 42.2 \approx
\qty{930}{kJ}$ épargnés sur la course, environ 9~\%.

\textbf{19.} \omterm{def:g10:exercise-and-energy:glycogen}{Glycogène} \qty{8500}{kJ}~; \omterm{def:g10:chemistry-of-life:families}{glucides} nécessaires $0.7 \times
8360 \approx \qty{5850}{kJ}$~: oui, les réserves suffisent sans boire.

\textbf{20.} Sans préparation, le \omterm{def:g10:exercise-and-energy:glycogen}{glycogène} de Nadia s'épuise vers le
kilomètre 34~; la surcharge glucidique avant la course, la boisson
glucidique pendant, et un entraînement qui élève la part de graisse
brûlée repoussent chacun le mur au-delà de la ligne d'arrivée.
\end{solution}
